
\documentclass[10pt]{article} % For LaTeX2e
\usepackage{tmlr}
% If accepted, instead use the following line for the camera-ready submission:
%\usepackage[accepted]{tmlr}
% To de-anonymize and remove mentions to TMLR (for example for posting to preprint servers), instead use the following:
%\usepackage[preprint]{tmlr}

% Optional math commands from https://github.com/goodfeli/dlbook_notation.
\input{math_commands.tex}

\usepackage{hyperref}
\usepackage{url}
\usepackage{graphicx}
\usepackage{booktabs}
\usepackage{multirow}
\usepackage{amsmath}

\title{Compression Favors Consistency, Not Truth: When and Why Language Models Prefer Correct Information}

% Authors must not appear in the submitted version. They should be hidden
% as long as the tmlr package is used without the [accepted] or [preprint] options.
% Non-anonymous submissions will be rejected without review.

\author{\name Konstantin Krestnikov \email krestnikov@example.com \\
      \addr Independent Researcher}

\newcommand{\fix}{\marginpar{FIX}}
\newcommand{\new}{\marginpar{NEW}}

\def\month{MM}  % Insert correct month for camera-ready version
\def\year{YYYY} % Insert correct year for camera-ready version
\def\openreview{\url{https://openreview.net/forum?id=XXXX}} % Insert correct link to OpenReview for camera-ready version


\begin{document}


\maketitle

\begin{abstract}
Why do language models trained on contradictory data prefer correct answers? In controlled experiments with small transformers (3.5M--86M parameters), we show that this preference tracks the compressibility structure of errors rather than truth per se. We train GPT-2 style models on corpora where each mathematical problem appears with both correct and incorrect solutions -- a denoising design that directly models conflicting information about the same fact. When errors are random, models extract the correct signal with accuracy scaling from 65\% to 85\% with model size. When errors follow a coherent alternative rule system, accuracy drops to chance (${\sim}$45--51\%): the model cannot distinguish the false system from truth. A multi-rule experiment reveals a sharp crossover: a single coherent alternative rule eliminates truth bias entirely, but adding a second competing rule restores most of it (47\%$\to$78\%), with continued growth through $N{=}10$ (88\%). The same pattern reproduces on real Wikipedia text (71\% vs 46\%). We propose the Compression--Consistency Principle as an explanatory hypothesis: in these settings, gradient descent favors the most compressible answer cluster, not truth per se. Truth bias emerges only when falsehood is structurally incoherent. Whether this principle extends to large-scale pretraining remains an open question.
\end{abstract}

\section{Introduction}
\label{sec:intro}

Language models train on data that contains contradictions. The internet mixes correct and incorrect claims about the same facts, often in the same format and style. When a model compresses such a corpus, which version does it favor?

Several explanations have been proposed. Scaling helps: larger models perform better on factual tasks \citep{kadavath2022language}. RLHF steers models toward human-preferred outputs. Data statistics play a role: factual accuracy correlates with frequency and source reliability \citep{elazar2022measuring,joshi2024personas,kandpal2023large}. Internal truth representations have been discovered in model activations \citep{burns2023discovering,marks2023geometry}. Yet none of these address a more fundamental question: \textit{why would the training objective itself -- next-token prediction -- create any preference for truth?}

We propose that the answer lies in compression. Minimizing cross-entropy is equivalent to minimizing code length \citep{shannon1948mathematical,deletang2024language}, connecting LLM training to the Minimum Description Length principle \citep{rissanen1978modeling,grunwald2007minimum}. But compression does not inherently favor truth -- it favors the most \textit{compressible} hypothesis consistent with the data.

We propose the \textbf{Compression--Consistency Principle} as an explanatory hypothesis:

\begin{quote}
In our settings, models do not privilege truth directly; they privilege the most compressible answer structure available in the data.
\end{quote}

Truth benefits from compression only when falsehood is structurally incoherent. Diverse errors must be memorized individually, whereas a correct rule system compresses into a compact representation. When errors form a coherent alternative system -- internally consistent but wrong -- they compress just as efficiently, and the preference vanishes.

We test this with a \textbf{denoising design}: each mathematical problem appears in the training corpus with both correct and incorrect solutions. This directly models the scenario where the same question receives conflicting answers. Four conditions vary the structure and ratio of contradictory answers (Section~\ref{sec:methodology}). The results are unambiguous: random errors produce strong, scaling truth bias; coherent errors produce none. Multi-rule experiments reveal a sharp crossover at the boundary (Section~\ref{sec:multirule}). Wikipedia entity substitution demonstrates the same pattern on real text (Section~\ref{sec:wikipedia}). Supporting experiments -- standard baselines, compression measurement, robustness checks, and verification dependencies -- are reported in the appendices.

Three caveats apply throughout. First, models compress \textit{text}, not reality: ``truth'' here means correctness of mathematical derivations and factual accuracy of Wikipedia paragraphs. Second, frequency can override compressibility -- when errors dominate sufficiently, the model follows the majority. Third, the compressibility gap is corpus-dependent: in natural language, it is smaller than in formal math.

The work makes five contributions:

\begin{enumerate}
    \item A \textbf{denoising experimental design} where the same problem appears with contradictory answers, directly testing signal extraction from noise.
    \item The \textbf{random/coherent contrast} as the central finding: random errors yield strong truth bias that scales with model capacity; coherent errors yield none.
    \item A \textbf{noise tolerance curve} showing graceful degradation as the signal-to-noise ratio decreases, with capacity-dependent plateaus.
    \item A \textbf{multi-rule sharp crossover}: a single coherent alternative rule eliminates truth bias; two competing rules restore most of it (47\%$\to$78\%), with accuracy continuing to grow through $N{=}10$ (88\%).
    \item \textbf{Transfer to real text}: a Wikipedia entity-substitution experiment reproduces the same random/coherent contrast on natural language.
\end{enumerate}

\section{Related Work}
\label{sec:related}

\subsection{Prediction as Compression}

The link between prediction and compression traces back to \citet{shannon1948mathematical}, who showed that optimal compression requires knowledge of the true data distribution. \citet{solomonoff1964formal} formalized optimal prediction as weighting hypotheses by program length. \citet{rissanen1978modeling} developed the Minimum Description Length (MDL) principle: the best model minimizes total description length. \citet{grunwald2007minimum} systematized MDL and showed its equivalence to several forms of inference. \citet{hutter2005universal} developed these ideas into a formal theory of universal intelligence (AIXI), linking intelligence to compression.

In the context of language models, \citet{deletang2024language} showed that LLMs are universal compressors. \citet{huang2024compression} discovered a linear correlation ($r \sim -0.95$) between compression quality and benchmark performance. \citet{wan2025large} proved that LLM training approximates Solomonoff induction. \citet{pan2025understanding} used compression-based analysis to explain knowledge acquisition and scaling. \citet{chlon2025predictable} linked compression failures to hallucination patterns. Our work builds directly on the MDL framework: we experimentally vary the description length of false answer systems and observe under what conditions the MDL-optimal choice coincides with truth.

\subsection{Internal Representations of Truth in LLMs}

Several studies have found internal representations correlated with truthfulness. \citet{marks2023geometry} showed a linear geometric structure of truthfulness in activation space. \citet{burns2023discovering} proposed CCS for discovering truth directions without supervision. \citet{li2023b_inference} identified a gap between internal knowledge and generation. \citet{ravfogel2025emergence} proposed the Truth Co-occurrence Hypothesis. \citet{azaria2023internal} showed that internal states can distinguish true from false outputs. \citet{burger2024truth} demonstrated robust cross-model lie detection. \citet{halawi2024overthinking} found that larger models can ``overthink'' and revert to memorized facts. \citet{ortu2024competition} traced competition between factual recall and counterfactual signals. Our work complements these at the behavioral level: we study when compression produces a preference for correct completions, leaving activation-level analysis for future work.

\subsection{Emergent World Models}

Language models can form internal world models from pure text prediction. \citet{li2023a_emergent} found that a model trained on Othello moves learns a board representation. \citet{gurnee2024language} discovered linear representations of space and time in Llama-2. These show that compression can give rise to structured representations. Our work asks when compression yields behavioral preference for correct versus incorrect completions.

\subsection{Truthfulness and Training Data Statistics}

\citet{joshi2024personas} showed that truthfulness links to ``personas'' in pretraining data. \citet{elazar2022measuring} demonstrated dependence on fact frequency. \citet{kandpal2023large} showed a direct relationship between supporting documents and accuracy. \citet{kalai2024calibrated} proved that calibrated models must hallucinate at a rate tied to the corpus's monofact rate. Our work differs: we \textit{experimentally vary} error structure while controlling frequency, isolating compressibility as a factor beyond frequency and source reliability.

\subsection{Simplicity Bias, Noisy Labels, and Grokking}

Neural networks prefer simple functions \citep{valle2019deep,mingard2021sgd,goldblum2024no,bhattamishra2023simplicity}. \citet{mesaros2024rule} showed that models generalize compositionally when rules are simple. The noisy labels literature directly parallels our setup: \citet{zhang2017understanding} showed that networks can memorize random labels but generalize through structure; \citet{rolnick2017deep} showed robustness to massive label noise. Grokking \citep{nanda2023progress,demoss2024complexity,liu2023grokking} also connects to compression: the transition from memorization to generalization is a compression process. Our experiments with coherent errors show that a coherent false system is just as ``simple'' as truth, and compression shows no preference.

\subsection{Our Contribution}

The works above study internal representations, establish theoretical links between compression and intelligence, or analyze dependence on data statistics. To our knowledge, direct training experiments with systematic variation of \textit{error compressibility} in a \textit{denoising} setting -- where contradictory answers coexist for the same problem -- remain limited. This work contributes such a controlled study. Unlike frequency-based analyses, we fix frequency and vary error structure. Unlike representational studies, we train from scratch and identify behavioral conditions for preference. Unlike noisy-label work, we show that ``structured noise'' (coherent errors) is not filtered out.

\section{Methodology}
\label{sec:methodology}

\subsection{Model and Training}
\label{sec:model}

GPT-2 style decoder-only transformers with pre-norm (LayerNorm before attention/MLP), GELU activation, and causal mask.

\begin{table}[t]
\caption{Model configurations.}
\label{tab:configs}
\begin{center}
\begin{tabular}{lcccc}
\toprule
Config & Layers & $d_{\text{model}}$ & Heads & Parameters \\
\midrule
tiny & 4 & 256 & 4 & ${\sim}$3.5M \\
small & 6 & 384 & 6 & ${\sim}$12M \\
medium & 8 & 512 & 8 & ${\sim}$26M \\
large & 12 & 768 & 12 & ${\sim}$86M \\
\bottomrule
\end{tabular}
\end{center}
\end{table}

Denoising experiments (Section~\ref{sec:j1j2}--\ref{sec:j3j4}) use a PyTorch implementation; multi-rule (Section~\ref{sec:multirule}), Wikipedia (Section~\ref{sec:wikipedia}), and standard experiments (Appendix~\ref{app:standard}) use an MLX implementation. Exact parameter counts vary slightly by tokenizer and framework (within 10\%). Optimizer: AdamW (weight\_decay=0.01), cosine decay with linear warmup, lr=$3 \times 10^{-4}$, seq\_len=256, batch\_size=32, 5000 steps. All experiments are repeated with multiple random initializations (4 seeds for core conditions, 2 seeds where noted). Learning curves show behavioral plateau by step 3000--4000 at all sizes (Appendix~\ref{app:robustness}).

\subsection{Denoising Corpus Design}
\label{sec:corpus}

The denoising setup is the primary experimental paradigm. The generator creates mathematical problems of four types: multi-step arithmetic, factorization, equation solving, and differentiation. Each problem is formatted as a step-by-step derivation in English, verified by SymPy. The tokenizer is character-level (vocab ${\sim}$57). Each denoising corpus contains 5,000 unique problems (10,000+ texts for J1/J2, 15,000 for J3, 25,000 for J4). Standard corpora (Appendix~\ref{app:standard}) contain ${\sim}$200,000 problems (${\sim}$36MB). Paired evaluation uses ${\sim}$5,000 held-out test pairs per condition, generated with a separate seed.

The key design feature: \textbf{each problem appears multiple times in the corpus with contradictory answers.} This directly models the internet scenario where the same question receives conflicting responses.

\textbf{Example (J1, random errors):}
\begin{verbatim}
Problem: Simplify 2 * (x + 3)
Correct:  2 * (x + 3) = 2x + 6      [appears once]
Wrong:    2 * (x + 3) = 2x + 5      [appears once, random error]
\end{verbatim}

\textbf{Example (J2, coherent errors):}
\begin{verbatim}
Problem: Simplify 2 * (x + 3)
Correct:  2 * (x + 3) = 2x + 6      [appears once]
Wrong:    2 * (x + 3) = 2 * (x + 2)
          = 2x + 4                    [appears once, systematic rule:
                                       a*(b+c) = a*(b+c-1)]
\end{verbatim}

Four conditions vary the structure and ratio of contradictory answers:

\begin{table}[t]
\caption{Denoising conditions.}
\label{tab:conditions}
\begin{center}
\begin{tabular}{lcccl}
\toprule
Condition & Correct & Incorrect & Ratio & Purpose \\
\midrule
\textbf{J1} & 1 & 1 random & 1:1 & Core: signal extraction from random noise \\
\textbf{J2} & 1 & 1 coherent & 1:1 & Control: coherent errors eliminate bias \\
\textbf{J3} & 1 & 2 random & 1:2 & Noise tolerance at moderate noise \\
\textbf{J4} & 1 & 4 random & 1:4 & Noise tolerance at high noise \\
\bottomrule
\end{tabular}
\end{center}
\end{table}

A fifth exploratory condition (J5: zero correct answers) is reported in Appendix~\ref{app:j5}.

\subsection{Wikipedia Corpus Design}
\label{sec:wiki_corpus}

To test transfer beyond formal math, we construct a real-text corpus from Wikipedia articles. Using NER (spaCy), we identify named entities in Wikipedia paragraphs and create two corruption modes:

\begin{itemize}
    \item \textbf{Random substitution:} Each entity is replaced with a random entity of the same type from a global pool (e.g., ``Paris'' $\to$ ``Kumamoto'').
    \item \textbf{Coherent substitution:} A consistent mapping is applied across the entire corpus (e.g., every ``France'' $\to$ ``Japan'', every ``Paris'' $\to$ ``Tokyo''), preserving internal consistency.
\end{itemize}

The corpus contains 20,000 paragraphs from 10,000 Wikipedia articles. Models are trained on a 50/50 mix of original and corrupted paragraphs. Evaluation uses paired NLL comparison on 2,000 held-out paragraphs: original vs.\ corrupted continuation given a shared context prefix.

\subsection{Evaluation Metrics}
\label{sec:eval}

\textbf{Paired evaluation (primary metric).} For each problem, a shared prompt is generated along with two completions (correct and incorrect). NLL is computed only on completion tokens, conditioned on the shared prompt. Metrics: \textbf{pair accuracy} (fraction of pairs where the model assigns lower NLL to the correct completion; our primary metric), mean $\Delta$Loss on completions, and Wilcoxon signed-rank test. Pair accuracy is equivalent to the Common Language Effect Size (CLES; \citealp{mcgraw1992common}).

\textbf{Corpus-level evaluation (secondary diagnostic).} We report corpus-level $\Delta$Loss $=$ Loss(incorrect) $-$ Loss(correct) as a secondary check. This metric is sensitive to text statistics confounds (Appendix~\ref{app:standard}) and should not be treated as the primary measure of truth bias.

\textbf{Statistical analysis.} For training variability we report seed-level dispersion. For paired evaluation, the one-sided Wilcoxon test quantifies uncertainty over held-out pairs within a single trained model. With thousands of pairs, $p$-values are inevitably minuscule; readers should focus on pair accuracy and seed-level dispersion rather than $p$-values.

\subsection{MDL Heuristic Framing}
\label{sec:mdl}

We state a heuristic MDL interpretation \citep{rissanen1978modeling,grunwald2007minimum}. Consider a denoising corpus where each problem appears with a correct answer (theory $T_1$) and an alternative answer (theory $T_2$). An idealized MDL learner minimizes $L(M) + L(D|M)$.

\textbf{Heuristic predictions:}

\begin{enumerate}
    \item $K(T_2) \gg K(T_1)$ \textbf{(random errors).} Each false answer requires independent encoding. The false system's description length grows with corpus size. An MDL learner should favor $T_1$.

    \item $K(T_2) \approx K(T_1)$ \textbf{(coherent errors).} Both systems are described by compact rules of comparable complexity. At equal frequency, an MDL learner has no reason to prefer one over the other.

    \item \textbf{Increasing noise ratio.} As the number of random false answers per problem grows (J3, J4), the correct signal becomes a smaller fraction, but each individual false answer remains unique, preserving the compressibility advantage of truth -- up to a capacity-dependent limit.

    \item \textbf{Multiple alternative rules.} $N$ competing false rules are individually compact but collectively incoherent: the mapping ``problem $\to$ which rule'' is unpredictable. Total description length of the false cluster grows with $N$, restoring the MDL advantage of truth.
\end{enumerate}

The denoising experiments directly test these predictions.

\section{Results}
\label{sec:results}

\subsection{J1 vs J2: The Central Contrast}
\label{sec:j1j2}

\begin{table}[t]
\caption{Denoising paired evaluation: J1 (1 correct + 1 random) vs J2 (1 correct + 1 coherent), across model sizes.}
\label{tab:j1j2}
\begin{center}
\begin{tabular}{lccccccc}
\toprule
Size & Params & J1 Acc. & J1 $\Delta$Loss & J2 Acc. & J2 $\Delta$Loss & J1 Seeds & J2 Seeds \\
\midrule
tiny & 3.5M & \textbf{65.3\% $\pm$ 1.3\%} & +0.015 & 43.5\% $\pm$ 2.6\% & $-$0.007 & 4 & 4 \\
small & 12M & \textbf{74.6\% $\pm$ 1.6\%} & +0.027 & 44.5\% $\pm$ 3.0\% & $-$0.004 & 4 & 4 \\
medium & 26M & \textbf{81.1\% $\pm$ 1.2\%} & +0.037 & 45.8\% $\pm$ 3.4\% & $-$0.003 & 4 & 4 \\
large & 86M & \textbf{85.2\% $\pm$ 2.3\%} & +0.048 & 51.0\% $\pm$ 0.8\% & $-$0.001 & 2 & 2 \\
\bottomrule
\end{tabular}
\end{center}
\end{table}

\begin{figure}[t]
\begin{center}
\includegraphics[width=0.85\linewidth]{results/figure_denoising_j1_j2.png}
\end{center}
\caption{The central contrast. J1 (random errors): accuracy scales monotonically from 65\% to 85\%. J2 (coherent errors): accuracy stays near chance (44--51\%) across all sizes. The gap between the curves is the compression-consistency effect.}
\label{fig:j1j2}
\end{figure}

This is the central result. When the same problem appears with both a correct and a random wrong answer, the model progressively learns to prefer the correct one -- accuracy reaches 85\% at 86M parameters. When the wrong answer follows a coherent rule system, the effect disappears: accuracy hovers near 50\% at all scales.

\textbf{Scaling.} J1 accuracy increases monotonically: 65\% $\to$ 75\% $\to$ 81\% $\to$ 85\%. The largest gain is between tiny and small (+9 pp), with continued growth through large. J2 accuracy converges toward 50\% from below as models grow, consistent with the MDL prediction that equal-description-length systems at equal frequency should be indistinguishable.

\textbf{Below-chance accuracy at small scale (J2).} The tiny model's J2 accuracy of 43.5\% is consistently below 50\%, replicating across multiple independent setups: standard coherent (47.2\%), BPE standard coherent (45.9\%), BPE denoising coherent (49.3\%), synthetic world coherent (46.6\%). This reflects textual simplicity asymmetry: the coherent error rules produce shorter or simpler outputs than the true rules (e.g., dropping a coefficient in derivatives). As capacity grows, the model represents both systems with equal per-token accuracy, and the asymmetry disappears.

This below-chance pattern is not a confound -- it strengthens the thesis. When the false system is textually simpler, the compressor \textit{actively prefers it over truth}. This is exactly what the Compression--Consistency Principle predicts: the compressor has no concept of correctness; it follows whichever signal is easier to encode. When coherent errors are simpler, they win; when matched in complexity, the result is chance; in this controlled setting, truth receives no preferential treatment.

A potential concern is that the entire random/coherent contrast reduces to local text complexity rather than system-level compressibility. Three controls address this:

\begin{enumerate}
    \item \textbf{Multi-rule control.} The multi-rule experiment (Section~\ref{sec:multirule}) uses the \textit{same} error rules with the \textit{same} textual complexity at $N{=}1$ and $N{=}2$ -- only the number of competing rules changes. If local simplicity drove the effect, $N{=}2$ should behave like $N{=}1$. Instead, accuracy jumps from 47\% to 78\% at tiny and to 86\% at small, confirming that the critical variable is the diversity of error rules, not their surface form.

    \item \textbf{Length-matched evaluation.} Our eval code computes both standard mean NLL (per-token, averaged over completions) and length-matched mean NLL (controlling for completion length differences). In the standard setup, mean completion lengths are nearly identical: 81.2 tokens (correct) vs 80.7 tokens (incorrect) for random, 88.4 vs 87.6 for coherent. Length-matched accuracy is within 1 percentage point of standard accuracy across all conditions: 78.7\% vs 79.5\% (random tiny), 85.8\% vs 86.2\% (random small), 46.2\% vs 47.2\% (coherent tiny). The effect is not driven by length artifacts.

    \item \textbf{Sum-NLL metric.} We also report sum-NLL accuracy (total NLL on the completion, not per-token average), which penalizes longer completions. Sum-NLL accuracy matches per-token accuracy (e.g., 78.4\% vs 79.5\% for random tiny), further ruling out length confounds.
\end{enumerate}

\subsection{J3 and J4: Noise Tolerance}
\label{sec:j3j4}

What happens when the noise ratio increases? J3 adds two random wrong answers per problem (1:2 ratio); J4 adds four (1:4 ratio).

\begin{table}[t]
\caption{Noise tolerance: paired accuracy across denoising conditions.}
\label{tab:noise}
\begin{center}
\begin{tabular}{lcccccc}
\toprule
Condition & Ratio & Tiny & Small & Medium & Large & Seeds \\
\midrule
J1 & 1:1 & 65.3\% & 74.6\% & 81.1\% & 85.2\% & 4/4/4/2 \\
J3 & 1:2 & 59.0\% & 68.6\% & 73.6\% & 75.2\% & 2/2/2/2 \\
J4 & 1:4 & 56.6\% & 64.9\% & 66.6\% & 65.8\% & 2/2/2/2 \\
J2 & 1:1 coh. & 43.5\% & 44.5\% & 45.8\% & 51.0\% & 4/4/4/2 \\
\bottomrule
\end{tabular}
\end{center}
\end{table}

\begin{figure}[t]
\begin{center}
\includegraphics[width=0.85\linewidth]{results/figure_denoising_noise_curve.png}
\end{center}
\caption{Noise tolerance. Left: accuracy vs model size for J1/J3/J4. Higher noise ratios yield lower accuracy and earlier plateaus. Right: accuracy vs noise ratio at each model size, showing graceful degradation.}
\label{fig:noise}
\end{figure}

Three observations:

\begin{enumerate}
    \item \textbf{Graceful degradation.} Increasing the noise ratio reduces accuracy but does not destroy the signal. At 1:4 noise (J4), even the tiny model extracts a weak signal (56.6\%), and the small model reaches 65\%.

    \item \textbf{Capacity-dependent plateaus.} J4 accuracy plateaus at medium/large: 66.6\% $\to$ 65.8\%. The signal-to-noise ratio becomes the bottleneck -- additional capacity cannot overcome 4:1 noise. J3 still shows growth at large (75.2\%), though slower than J1. J1 shows no plateau in the tested range.

    \item \textbf{Signal-to-noise logic.} For each noise ratio, there exists a model capacity at which accuracy saturates. Lower ratios saturate at higher accuracy. This is consistent with the MDL framing: the compressibility advantage of the correct cluster persists but must compete with sheer volume of noise.
\end{enumerate}

\subsection{Multi-Rule Errors: The Sharp Crossover}
\label{sec:multirule}

The denoising experiments establish two poles: one coherent rule yields chance, random errors yield strong bias. What lies between? We train models on corpora with $N$ alternative wrong rules per task type; for each problem, one rule is chosen at random. Each rule is compact, but the mapping ``problem $\to$ rule'' is unpredictable.

\begin{table}[t]
\caption{Matched paired evaluation of multi-rule errors (tiny, 3.5M, 50/50, 4 seeds).}
\label{tab:multirule}
\begin{center}
\begin{tabular}{cccc}
\toprule
$N$ Rules & Accuracy & $\Delta$Loss (paired) & Wilcoxon $p$ \\
\midrule
1 (coherent) & 46.6\% & $-$0.0019 & ${\sim}$1.0 \\
2 & 77.6\% & +0.0152 & $< 10^{-6}$ \\
3 & 82.8\% & +0.0213 & $< 10^{-6}$ \\
5 & 84.8\% & +0.0293 & $< 10^{-6}$ \\
10 & 88.3\% & +0.0440 & $< 10^{-6}$ \\
\bottomrule
\end{tabular}
\end{center}
\end{table}

The sharp crossover reproduces at larger scale: $N{=}2$ on small (12M) yields 86.3\% $\pm$ 0.8\% (4 seeds), compared to 77.6\% on tiny. The effect strengthens with capacity, consistent with the MDL hypothesis.

\begin{figure}[t]
\begin{center}
\includegraphics[width=0.7\linewidth]{results/figure7_multirule.png}
\end{center}
\caption{Multi-rule sharp crossover. Accuracy jumps from 47\% at $N{=}1$ to 78\% at $N{=}2$, then continues growing gradually through $N{=}10$.}
\label{fig:multirule}
\end{figure}

The largest increase is between one rule and two rules (47\% $\to$ 78\%), but the curve remains graded afterward. This is the most mechanistically revealing result: a single coherent alternative rule suffices to eliminate truth bias entirely, while even two alternatives restore most of it.

The MDL hypothesis offers a natural explanation: with one false rule, $K(T_2) \approx K(T_1)$ -- both are compact. With two false rules, the learner must additionally encode a \textit{selector}: which rule applies to which problem. This selector is essentially a random function from problems to $\{\text{rule}_1, \text{rule}_2\}$ -- it has high Kolmogorov complexity because there is no pattern the compressor can exploit. The selector's description length cost is what breaks the false system's compressibility advantage, even though each individual rule remains compact. As $N$ grows, the selector becomes increasingly complex ($\log N$ bits per problem), and the false cluster approaches the incompressibility of fully random errors. At $N{=}10$, accuracy (88.3\%) exceeds even the standard random baseline (79.5\%).

\subsection{Transfer to Real Text: Wikipedia}
\label{sec:wikipedia}

Does the principle transfer beyond formal math? We train models on Wikipedia paragraphs where entities have been substituted (Section~\ref{sec:wiki_corpus}). Paired evaluation compares original paragraphs against corrupted versions given a shared context prefix.

\begin{table}[t]
\caption{Wikipedia entity substitution (50/50, 4 seeds per size).}
\label{tab:wiki}
\begin{center}
\begin{tabular}{lcccc}
\toprule
Size & Random Acc. & Random $\Delta$Loss & Coherent Acc. & Coherent $\Delta$Loss \\
\midrule
tiny (3.5M) & \textbf{69.6\% $\pm$ 0.1\%} & +0.064 & 48.7\% $\pm$ 0.5\% & +0.000 \\
small (12M) & \textbf{70.7\% $\pm$ 0.4\%} & +0.065 & 46.6\% $\pm$ 0.4\% & $-$0.002 \\
medium (26M) & \textbf{71.5\% $\pm$ 0.8\%} & +0.068 & 46.4\% $\pm$ 0.8\% & $-$0.003 \\
large (86M) & \textbf{71.4\% $\pm$ 0.7\%} & +0.069 & 45.9\% $\pm$ 1.5\% & $-$0.005 \\
\bottomrule
\end{tabular}
\end{center}
\end{table}

\begin{figure}[t]
\begin{center}
\includegraphics[width=0.85\linewidth]{results/figure_wiki_results.png}
\end{center}
\caption{Wikipedia entity substitution. Left: random accuracy is 70--71\% across all sizes, well above chance. Coherent accuracy is at or below chance (46--49\%). Right: per-entity-type accuracy for random errors (tiny, seed 42).}
\label{fig:wiki}
\end{figure}

The random/coherent contrast reproduces on real text:

\begin{enumerate}
    \item \textbf{Truth bias reproduces.} Random substitution yields 70--71\% accuracy across all sizes (all seeds $p < 10^{-100}$). The model distinguishes original Wikipedia paragraphs from randomly corrupted ones.

    \item \textbf{Coherent substitution eliminates the effect.} Coherent substitution yields 46--49\% accuracy -- at or below chance. When substitution is internally consistent (every ``France'' $\to$ ``Japan'', every ``Paris'' $\to$ ``Tokyo''), the model cannot distinguish original from corrupted text.

    \item \textbf{The effect is weaker than in math.} Wikipedia random accuracy (71\%) is substantially below math denoising J1 at comparable scales (65--85\%). Natural language provides more textual flexibility, allowing corrupted text to remain locally fluent.

    \item \textbf{Scaling is flat.} Unlike math, where accuracy grows from 65\% at tiny to 85\% at large, Wikipedia accuracy saturates near 70\% across the tested size range. The compressibility gap in natural language is smaller and does not substantially benefit from additional model capacity in the 3.5M--86M range.
\end{enumerate}

\textbf{Per-entity-type accuracy} (random, tiny):

\begin{table}[t]
\caption{Per-entity-type accuracy (Wikipedia, random, tiny).}
\label{tab:wiki_entity}
\begin{center}
\begin{tabular}{lcc}
\toprule
Entity Type & $N$ Pairs & Accuracy \\
\midrule
LOC & 45 & 82.2\% \\
NORP & 197 & 78.7\% \\
GPE & 301 & 77.1\% \\
PERSON & 402 & 69.9\% \\
DATE & 211 & 67.3\% \\
ORG & 664 & 65.2\% \\
CARDINAL & 172 & 61.0\% \\
\bottomrule
\end{tabular}
\end{center}
\end{table}

Geographic entities (GPE, LOC, NORP) show the strongest effect (77--82\%), likely because location names create strong contextual expectations. Organization and person names are more interchangeable. Cardinal numbers are weakest (61\%), consistent with their high variability in natural text.

\section{Discussion}
\label{sec:discussion}

\subsection{The Compression--Consistency Principle}

The experiments support a unified explanatory hypothesis:

\textbf{In our controlled settings, the compression objective tracks consistency rather than truth.} Any internally consistent rule system -- true or false -- compresses equally well. Truth bias emerges only when false alternatives are structurally incoherent. We frame this as the Compression--Consistency Principle -- a hypothesis that organizes our observations, not a proven mechanism of language model training in general.

The evidence forms a coherent picture across all experimental conditions:

\begin{enumerate}
    \item \textbf{Random errors are incompressible.} Each random error is unique; the false cluster cannot compress into a compact rule. The correct cluster has lower description length, and the model preferentially encodes it. This holds in the denoising setup (J1: 65--85\%), in the standard setup (80--88\%), and on Wikipedia (70--71\%).

    \item \textbf{Coherent errors are equally compressible.} One systematic wrong rule is as compact as one correct rule. Compression provides no basis to prefer truth. This holds in denoising (J2: 44--51\%), in the standard setup (47--52\%), on Wikipedia (46--49\%), and in a synthetic world domain (46.6\%).

    \item \textbf{Rule diversity is the key variable.} The multi-rule experiment provides the sharpest test: increasing the number of false rules from 1 to 10 progressively degrades the compressibility of the false cluster, producing a graded curve from chance to 88\%. The sharp crossover at $N{=}1 \to 2$ (47\%$\to$78\%) is the most informative data point: it takes only one additional rule to break the compressibility of the false cluster.

    \item \textbf{The pattern transfers across domains.} Math (formal derivations), Wikipedia (factual paragraphs), and a synthetic world (fictional entities with deterministic rules) all reproduce the random/coherent contrast. The effect size varies -- strongest in math, weakest in the synthetic world -- but the direction is consistent across all tested domains.

    \item \textbf{Compressibility correlates with accuracy.} A gzip compression analysis on raw completion text shows a positive rank correlation with paired accuracy across 9 conditions (Spearman $\rho = 0.68$, $p = 0.042$; Appendix~\ref{app:compression}). While the sample size is small, the direction is consistent: conditions with larger compression gaps produce stronger truth bias.

    \item \textbf{Frequency alone does not explain the results.} In the standard setup, truth bias persists even at 10/90 mixing (67\% paired accuracy; Appendix~\ref{app:standard}) -- the model prefers correct answers despite 9:1 frequency disadvantage. But for coherent errors, frequency is all that matters: at 40/60 coherent, the model follows the majority with 72\% preference for the false system (Appendix~\ref{app:standard}, Table~\ref{tab:coh_prop}). Compressibility overrides frequency for random errors; frequency governs for coherent ones.
\end{enumerate}

\subsection{Implications}

\textbf{For alignment.} In our controlled settings, the training objective does not provide an inherent ``truth compass.'' Systematic falsehood remains competitive when internally coherent -- a single consistent wrong rule is as compressible as the truth. If this pattern holds at larger scales, it would suggest that relying on scale alone may not solve the problem of coherent misinformation. Verification dependencies can partially restore truth bias (Appendix~\ref{app:robustness}), but their effectiveness decreases with model size in our experiments, raising questions about scalability.

\textbf{For understanding hallucinations.} Our results complement the statistical hallucination bound of \citet{kalai2024calibrated}. In our experiments, coherent misconceptions remain attractive to the model because they compress well, independently of their rarity. If this generalizes, hallucinations that happen to be internally consistent may be particularly persistent and resistant to correction through scale alone.

\textbf{For data curation.} In our experiments, when a corpus contains contradictory answers to the same question, models favor the more compressible cluster. This is reassuring when errors are diverse (as in naturally occurring mistakes) but concerning when errors are systematic. By analogy, coordinated campaigns that maintain internal consistency -- always substituting the same false narrative, always citing the same fabricated sources -- would function as a compressible alternative rule system. Our results suggest (but do not directly demonstrate at realistic scale) that such coherent falsehood could be harder for the compressor to distinguish from truth than diverse organic errors. The noise tolerance results (J3, J4) provide quantitative guidance on the optimistic side: even at 4:1 noise, small models extract truth with 65\% accuracy, as long as the noise is structurally incoherent.

\subsection{Analogy with Popper's Falsifiability}

Our results admit an interpretive analogy with the falsifiability criterion \citep{popper1959logic}: a true theory compresses maximally, while a false theory needs corrections that increase description length. However, the analogy is limited: the model does not ``test'' theories, regular discrepancies are themselves compressible, and ad hoc corrections do not produce truth bias (Appendix~\ref{app:standard}). For a compressor, a discrepancy is merely another pattern.

\subsection{Broader Impact}

This paper raises the possibility that internally consistent misinformation may be harder for language models to distinguish from truth than diverse, uncoordinated errors. We emphasize that this finding is established in a controlled small-scale setting (3.5M--86M parameters, synthetic math tasks and simple Wikipedia substitutions) and does not directly demonstrate vulnerability in production language models. The extrapolation to real-world disinformation campaigns (Section~5.2) is a hypothesis motivated by our results, not an empirically validated claim. We believe transparent reporting of the conditions under which truth bias fails is more beneficial than concealment, as it can inform defensive measures in data curation and model evaluation.

\subsection{Limitations}

\textbf{Model scale.} Experiments use models from 3.5M to 86M parameters. The denoising results show clear scaling trends but do not establish whether these continue beyond 86M. For coherent errors, there is a stronger heuristic argument: if both systems have comparable description length and equal frequency, an MDL learner has no basis to prefer one over the other regardless of scale.

\textbf{Domain specificity.} Mathematics has an unusually crisp correct/incorrect distinction. The effect weakens in natural language: 71\% on Wikipedia vs 65--85\% in math. Natural language provides more degrees of freedom for locally fluent errors, reducing the compressibility gap.

\textbf{Seed counts.} Core denoising conditions (J1, J2) use 4 seeds at tiny/small/medium and 2 seeds at large. Noise tolerance conditions (J3, J4) use 2 seeds. This is sufficient to show directional stability but does not tightly estimate between-run variability for the 2-seed conditions.

\textbf{Discriminative scope.} The primary metric is paired evaluation (forced choice). Generation sanity checks (Appendix~\ref{app:robustness}) confirm the direction (30.5\% vs 20.8\%) but the gap is smaller. The full relationship between discriminative truth bias and generative truthfulness remains an open question.

\textbf{J5 baseline.} An exploratory condition where models train on only incorrect answers (Appendix~\ref{app:j5}) shows 54--60\% pair accuracy for correct completions despite never seeing them. This could reflect intrinsic structural regularities in correct mathematical derivations (simpler numbers, consistent sign patterns) or residual format/length artifacts. While the length-matched analysis (Section~\ref{sec:j1j2}) rules out gross length confounds for the main results, J5 suggests that a small component of the observed effect may stem from the inherent textual regularity of correct answers. Importantly, this potential confound applies equally to J1 and J2; it cannot explain the random/coherent contrast, which is the paper's central finding.

\subsection{Future Work}

Key open directions:

\begin{enumerate}
    \item \textbf{Large-scale replication.} The most important next step is extending the denoising design to models at 1B+ parameters with compute-matched training budgets. Our results show clear trends from 3.5M to 86M; whether the random/coherent contrast persists, strengthens, or qualitatively changes at frontier scale is the critical open question for the practical relevance of the Compression--Consistency Principle.
    \item \textbf{Verification density.} Increasing cross-domain checks per task to test whether denser verification compensates for the compressor's growing power.
    \item \textbf{Real-world domains.} Extending to domains with competing knowledge systems (e.g., historical scientific debates) where the correct/incorrect boundary is well-established.
    \item \textbf{Internal representations.} Linear probing for ``truth directions'' vs ``coherence directions'' \citep{marks2023geometry} in models trained under our conditions.
    \item \textbf{Interaction with RLHF.} Testing whether reinforcement learning from human feedback amplifies or attenuates the compression-consistency effect.
    \item \textbf{Stronger surface-form controls.} While length-matched evaluation rules out gross length confounds and the multi-rule experiment controls for textual complexity, generating complexity-matched correct and incorrect completions would provide an even stronger test. The J5 baseline (Appendix~\ref{app:j5}) suggests that correct derivations may have intrinsic regularity advantages beyond what current controls capture.
\end{enumerate}

\section{Conclusion}
\label{sec:conclusion}

In controlled experiments with small transformers (3.5M--86M parameters), the compression objective tracks consistency rather than truth. When models train on contradictory answers to the same problems, they prefer the correct answer only when errors are structurally incoherent. In denoising experiments, accuracy scales from 65\% to 85\% for random errors, but stays near chance for coherent errors across all model sizes. The sharp crossover at $N{=}1 \to 2$ pinpoints the boundary: one coherent alternative rule eliminates truth bias entirely, two rules restore most of it. The same pattern reproduces on real Wikipedia text (71\% vs 46\%). Length-matched evaluation and the multi-rule complexity control rule out surface-form confounds as the primary driver. In the settings we study, truth bias is a compression artifact -- it emerges only when falsehood is structurally incoherent. Whether this principle extends to large-scale pretraining is the key open question; replication at 1B+ parameters is the natural next step.

\bibliography{references}
\bibliographystyle{tmlr}

\appendix

\section{Reproducibility}
\label{app:reproducibility}

All code, data generation scripts, and evaluation scripts will be made publicly available upon acceptance. Denoising experiments (Experiment J) were run on Modal.com T4 GPUs using PyTorch. Standard math experiments and Wikipedia experiments were run on an Apple Mac M4 with 36GB unified memory using the MLX framework (v0.31.0). Total computational cost: approximately 80 hours across all experiments. Over 160 models were trained across all conditions.

\section{Compression Measure}
\label{app:compression}

To operationalize the MDL argument, we measure the compression ratio of correct vs incorrect completion segments using gzip (level 9) on concatenated completions from each paired test set (thousands of completions per condition). The compression ratio delta provides a model-free proxy for the compressibility gap.

\begin{table}[t]
\caption{Compression ratio (gzip) and paired accuracy across conditions.}
\label{tab:compression}
\begin{center}
\begin{tabular}{lcccc}
\toprule
Condition & Correct & Incorrect & Delta & Paired Accuracy \\
\midrule
random 50/50 & 0.1627 & 0.1639 & +0.0012 & 79.5\% \\
coherent 50/50 & 0.1656 & 0.1658 & +0.0002 & 47.2\% \\
contradictory & 0.1745 & 0.1744 & $-$0.0001 & 49.0\% \\
multirule $N{=}2$ & 0.1671 & 0.1710 & +0.0038 & 77.6\% \\
multirule $N{=}3$ & 0.1669 & 0.1722 & +0.0053 & 82.8\% \\
multirule $N{=}5$ & 0.1672 & 0.1730 & +0.0057 & 84.8\% \\
multirule $N{=}10$ & 0.1676 & 0.1752 & +0.0076 & 88.3\% \\
world random & 0.0370 & 0.0516 & +0.0146 & 57.7\% \\
world coherent & 0.0374 & 0.0384 & +0.0010 & 46.6\% \\
\bottomrule
\end{tabular}
\end{center}
\end{table}

\begin{figure}[t]
\begin{center}
\includegraphics[width=0.7\linewidth]{results/figure_compression_vs_accuracy.png}
\end{center}
\caption{Compression ratio delta (gzip, incorrect minus correct) vs paired accuracy. Spearman $\rho = 0.68$, $p = 0.042$. Coherent and contradictory errors cluster near zero delta and chance accuracy; random and multi-rule errors show progressively larger deltas.}
\label{fig:compression}
\end{figure}

The rank correlation (Spearman $\rho = 0.68$, $p = 0.042$) is consistent with the hypothesis that compressibility gap drives truth bias, though the small sample size ($N{=}9$ conditions) limits the strength of this quantitative claim. Within the math domain, the multi-rule series is monotonic: larger compression deltas map to higher accuracy. We present this as supporting evidence, not as the primary basis for the compression-consistency principle, which rests on the experimental contrasts in Section~\ref{sec:results}.

\section{Standard (Non-Denoising) Experiments}
\label{app:standard}

\subsection{Standard Corpus Design}

As a complementary paradigm, we also train models on standard corpora where each problem appears once with either a correct or incorrect solution. The two groups do not share prompts.

\textbf{Error types:}
\begin{itemize}
    \item \textbf{Random:} A unique plausible error injected at a random step. Each error is different.
    \item \textbf{Coherent:} One systematic wrong rule per problem type. All problems of one type fail identically.
\end{itemize}

\subsection{Comparison with Denoising}

\begin{table}[t]
\caption{Standard vs denoising accuracy (50/50 random errors).}
\label{tab:std_vs_denoise}
\begin{center}
\begin{tabular}{lccc}
\toprule
Size & Standard & Denoising (J1) & Gap \\
\midrule
tiny & 79.5\% & 65.3\% & $-$14.2 pp \\
small & 86.2\% & 74.6\% & $-$11.6 pp \\
medium & 87.1\% & 81.1\% & $-$6.0 pp \\
large & 88.3\% & 85.2\% & $-$3.1 pp \\
\bottomrule
\end{tabular}
\end{center}
\end{table}

The denoising setup is substantially harder: within-problem contradiction is more challenging than across-corpus mixing. However, the gap closes with scale: 14 pp at tiny, 3 pp at large. Both setups show the same qualitative pattern.

\subsection{Standard Paired Baselines}

\begin{table}[t]
\caption{Standard paired evaluation at 50/50 (tiny, 4 seeds).}
\label{tab:std_baselines}
\begin{center}
\begin{tabular}{lccc}
\toprule
Error Type & Pair Accuracy & $\Delta$Loss (paired) & Wilcoxon $p$ \\
\midrule
Random & \textbf{79.5\%} & +0.032 & $< 10^{-6}$ \\
Contradictory & 49.0\% & +0.0003 & $> 0.3$ \\
Coherent & 47.2\% & $-$0.0018 & ${\sim}$1.0 \\
\bottomrule
\end{tabular}
\end{center}
\end{table}

\begin{table}[t]
\caption{Standard paired evaluation across proportions (random errors, 4 seeds).}
\label{tab:rand_prop}
\begin{center}
\begin{tabular}{ccc}
\toprule
Proportion & Pair Accuracy & $\Delta$Loss (paired) \\
\midrule
50/50 & 80\% & +0.032 \\
40/60 & 79\% & +0.043 \\
30/70 & 75\% & +0.036 \\
20/80 & 69\% & +0.029 \\
10/90 & 67\% & +0.017 \\
\bottomrule
\end{tabular}
\end{center}
\end{table}

Even at 10/90, where correct examples are outnumbered 9:1, paired accuracy remains 67\% ($p < 10^{-88}$). Truth bias persists under extreme frequency imbalance because the structural incompressibility of random errors provides a compression-based advantage that frequency alone cannot eliminate.

For coherent errors, the picture inverts: the model follows pure frequency, preferring whichever system is more abundant.

\begin{table}[t]
\caption{Random vs coherent paired accuracy across proportions.}
\label{tab:coh_prop}
\begin{center}
\begin{tabular}{ccc}
\toprule
Proportion & Random Accuracy & Coherent Accuracy \\
\midrule
50/50 & 80\% & 47.2\% \\
40/60 & 79\% & 27.8\% \\
30/70 & 75\% & 14.7\% \\
20/80 & 69\% & 9.6\% \\
\bottomrule
\end{tabular}
\end{center}
\end{table}

\begin{table}[t]
\caption{Standard paired accuracy by model size (50/50, 4 seeds).}
\label{tab:std_scaling}
\begin{center}
\begin{tabular}{lcc}
\toprule
Size & Random Accuracy & Coherent Accuracy \\
\midrule
tiny (3.5M) & 79.5\% & 47.2\% \\
small (12M) & 86.2\% & 49.6\% \\
medium (26M) & 87.1\% & 52.6\% \\
large (86M) & 88.3\% & 51.4\% \\
\bottomrule
\end{tabular}
\end{center}
\end{table}

\begin{figure}[t]
\begin{center}
\includegraphics[width=0.7\linewidth]{results/figure6_scaling.png}
\end{center}
\caption{Scaling in the standard setup. Random accuracy rises from 80\% to 88\%; coherent accuracy stays near chance (47--53\%) across the 3.5M--86M range.}
\label{fig:scaling}
\end{figure}

\subsection{Corpus-Level Evaluation and Its Limitations}

In the standard setup, corpus-level $\Delta$Loss provides a coarse measure. At 50/50 random, corpus-level $\Delta$Loss = +0.0115. The effect decreases monotonically with the fraction of correct data: +0.0089 (40/60), +0.0064 (30/70), +0.0033 (20/80), $-$0.0016 (10/90). The inversion at 10/90 is a frequency artifact: paired evaluation shows continued truth bias (67\%).

\subsection{Observations and Ad Hoc Correction (Experiments 2--3)}

We tested whether adding empirical feedback (observations showing discrepancies between false theory predictions and reality) restores truth bias for coherent errors.

\textbf{Observations (Experiment 2).} Adding 10--50\% observation checks to the coherent 50/50 corpus yields $\Delta$Loss = +0.0002 to +0.0008 -- an order of magnitude weaker than random errors. Regular discrepancies (the $a \times b = a \times (b-1)$ rule always understates by $a$) are themselves compressible.

\textbf{Five conditions for the false theory (Experiment 3).} Conditions A (no observations), B (bare discrepancies), C (ad hoc correction), D (systematic correction), E (non-specific predictions). Corpus-level $\Delta$Loss ordering: D $\approx$ C $>$ E $>$ B $\approx$ A. However, \textbf{paired evaluation shows no truth bias for any condition} (accuracy ${\sim}$49\% for all). The corpus-level effect is an artifact of text statistics differences between correct and incorrect corpora.

\subsection{NLL Distribution Analysis}

The per-pair NLL difference distribution explains metric divergence:
\begin{itemize}
    \item \textbf{Random errors:} right-skewed, mean = +0.048, median = +0.025, 81.5\% positive, heavy tail to +0.4.
    \item \textbf{Coherent errors:} symmetric around zero, 45.5\% positive.
\end{itemize}

The skew explains why pair accuracy and mean $\Delta$Loss can diverge: a minority of pairs with extreme $\Delta$Loss can pull the mean without affecting the median or the accuracy count.

\section{Additional Transfer Experiments}
\label{app:transfer}

\subsection{Synthetic World (Full Results)}

A synthetic world with 50 entities, 4 types, and 15 deterministic rules. Random 50/50: accuracy 57.7\% $\pm$ 1.7\%, $\Delta$Loss +0.034. Coherent 50/50: accuracy 46.6\% $\pm$ 1.7\%, $\Delta$Loss +0.019 (mixed -- the two metrics disagree in sign, likely due to distributional asymmetry). Per-type: minerals best (68.7\%), plants (62.2\%), animals near chance (51.4\%), potions near chance (49.1\%).

\subsection{Multi-Alternative Errors in the Synthetic World}

Multi-alternative errors ($N{=}1$ to $N{=}16$) produce a gradual rise in natural language: 47\% $\to$ 40\% $\to$ 50\% $\to$ 51\% $\to$ 60\%. Unlike math (47\% $\to$ 78\% at $N{=}2$), there is no steep early jump. At $N{=}2$, accuracy \textit{worsens} to 40\%: two alternatives collectively compete with truth while remaining compressible. Natural language absorbs contradictions that formal math cannot. This contrast between math and natural language is itself informative: it shows that the sharpness of the sharp crossover depends on how much redundancy the domain provides.

\subsection{Cross-Domain Falsification}

Adding correct cross-domain tasks (using true differentiation rules) to a coherent 50/50 corpus selectively increases derivative accuracy: 35\% $\to$ 56\% at 25\% cross-domain. The effect is non-monotonic (drops to 45\% at 50\% due to dilution). Other task types are unaffected. This provides preliminary evidence that cross-domain data can selectively weaken coherent errors when the correct rule is observed in a related context.

\section{Boundary Conditions and Robustness}
\label{app:robustness}

\subsection{Chained Verification}

If coherent errors cannot be distinguished from truth under compression alone, can embedding verification dependencies help? We construct \textit{chained tasks} where a computation using the coherent false rule is accompanied by arithmetic verification: for correct solutions, verification confirms the result (residual = 0); for coherent-error solutions, it produces an unpredictable numerical residual.

Six chain types: arithmetic $\to$ reverse, factoring $\to$ evaluation, linear equation $\to$ back-substitution, quadratic $\to$ root substitution, derivative $\to$ finite difference, tangent $\to$ prediction.

\begin{table}[t]
\caption{Chained tasks (50/50, coherent errors with verification).}
\label{tab:chained}
\begin{center}
\begin{tabular}{lccc}
\toprule
Size & Chained Accuracy & Standard Coherent & Improvement \\
\midrule
tiny (3.5M) & \textbf{70.9\% $\pm$ 1.2\%} & 43.3\% & +27.6 pp \\
small (12M) & 64.2\% $\pm$ 1.5\% & 49.6\% & +14.6 pp \\
large (86M) & 60.6\% $\pm$ 1.2\% & 51.4\% & +9.2 pp \\
\bottomrule
\end{tabular}
\end{center}
\end{table}

\begin{figure}[t]
\begin{center}
\includegraphics[width=0.7\linewidth]{results/figure10_chained.png}
\end{center}
\caption{Chained verification restores truth bias from 43\% to 71\% at tiny, but the effect weakens at larger model sizes. Control: truncated chains without verification yield 44.3\% -- same as standard coherent.}
\label{fig:chained}
\end{figure}

Key observations:

\begin{enumerate}
    \item \textbf{Verification restores truth bias.} Chained accuracy of 70.9\% is well above chance and above standard coherent errors. Verification transforms the coherent error's predictable output into an unpredictable verification residual, breaking its compressibility advantage.

    \item \textbf{The effect declines with model size.} Unlike random errors (80\% $\to$ 88\%), chained accuracy decreases: 71\% $\to$ 64\% $\to$ 61\%. A plausible explanation: higher-capacity models learn the coherent within-domain pattern more readily than the weaker verification signal. This is a striking result -- larger models may be \textit{better} at absorbing coherent errors, making verification less effective.

    \item \textbf{Accuracy varies by chain type.} Arithmetic reverse: 95.8\% (strongest -- incorrect multiplication yields a fractional result on reverse). Tangent prediction: 37.5\% (below chance -- $O(h^2)$ approximation error masks the coherent rule error).

    \item \textbf{Verification is the active ingredient.} Truncated chains (same tasks, no verification step) yield 44.3\% -- indistinguishable from standard coherent (43.3\%).
\end{enumerate}

\subsection{BPE Tokenization}

To rule out character-level encoding artifacts, we repeat the key comparison with a SentencePiece BPE tokenizer (vocab 1000).

\begin{table}[t]
\caption{BPE vs char-level paired accuracy, standard setup (tiny, 4 seeds).}
\label{tab:bpe_std}
\begin{center}
\begin{tabular}{lcc}
\toprule
Tokenizer & Random 50/50 & Coherent 50/50 \\
\midrule
Char (vocab 57) & 79.5\% $\pm$ 1.4\% & 47.2\% $\pm$ 2.8\% \\
BPE (vocab 1000) & \textbf{85.6\% $\pm$ 0.2\%} & \textbf{45.9\% $\pm$ 1.6\%} \\
\bottomrule
\end{tabular}
\end{center}
\end{table}

\begin{table}[t]
\caption{BPE vs char-level paired accuracy, denoising setup (tiny, 4 seeds).}
\label{tab:bpe_denoise}
\begin{center}
\begin{tabular}{lcc}
\toprule
Tokenizer & J1 (random) & J2 (coherent) \\
\midrule
Char (vocab 57) & 65.3\% $\pm$ 1.3\% & 43.5\% $\pm$ 2.6\% \\
BPE (vocab 1000) & \textbf{75.9\% $\pm$ 1.5\%} & \textbf{49.3\% $\pm$ 3.1\%} \\
\bottomrule
\end{tabular}
\end{center}
\end{table}

The effect survives BPE in both setups. In the denoising regime -- where contradictory answers coexist for the same problem -- BPE strengthens truth bias (75.9\% vs 65.3\% for J1) while coherent errors remain at chance (49.3\% for J2). This directly confirms that the denoising result is not an artifact of character-level tokenization.

\subsection{Generation Sanity Check}

Paired evaluation is a forced-choice setting. To verify that the effect extends to generation, we run greedy decoding on 500 test prompts per model with automated SymPy verification.

\begin{table}[t]
\caption{Generative accuracy (tiny, 4 seeds, 500 problems each).}
\label{tab:generation}
\begin{center}
\begin{tabular}{lcc}
\toprule
Condition & Mean Accuracy & $p$ (paired $t$-test) \\
\midrule
Random-trained & 30.5\% $\pm$ 1.7\% & 0.013 \\
Coherent-trained & 20.8\% $\pm$ 3.6\% & -- \\
\bottomrule
\end{tabular}
\end{center}
\end{table}

Random-trained models generate correct solutions 1.5$\times$ more often than coherent-trained ones. The generation gap is smaller than the paired gap (30\% vs 21\% generative, compared with 83\% vs 47\% discriminative), confirming that generation is harder. But the directional result holds.

\subsection{Learning Curves}

\begin{figure}[t]
\begin{center}
\includegraphics[width=0.85\linewidth]{results/figure_learning_curves.png}
\end{center}
\caption{Learning curves. All sizes reach behavioral plateau by step 3000--4000. The large model achieves 88.8\% at step 3000, stable through 5000. Coherent models stay near chance throughout.}
\label{fig:learning_curves}
\end{figure}

All sizes reach a behavioral plateau by step 3000--4000. This mitigates the concern about differential convergence across model sizes.

\section{J5 Zero-Signal Baseline (Exploratory)}
\label{app:j5}

J5 is the most provocative condition: the model trains on 2 random wrong answers per problem with no correct answers whatsoever. At evaluation, it is tested on paired comparison of correct vs random incorrect completions.

\begin{table}[t]
\caption{J5: zero-signal baseline (2 seeds per size).}
\label{tab:j5}
\begin{center}
\begin{tabular}{lccc}
\toprule
Size & Accuracy & $\Delta$Loss & Seeds \\
\midrule
tiny & 53.8\% $\pm$ 0.5\% & $-$0.046 & 2 \\
small & 59.9\% $\pm$ 1.0\% & $-$0.076 & 2 \\
medium & 58.1\% $\pm$ 2.7\% & $-$0.090 & 2 \\
large & 56.1\% $\pm$ 2.9\% & $-$0.111 & 2 \\
\bottomrule
\end{tabular}
\end{center}
\end{table}

The model assigns lower NLL to the correct answer in 54--60\% of pairs, despite never having seen correct answers during training. However, the mean $\Delta$Loss is strongly negative, indicating that on average the model assigns much lower NLL to the \textit{incorrect} answers it trained on. The distribution is left-skewed: most pairs slightly favor truth, but a minority strongly favor the trained-on errors.

\textbf{Interpretation.} One possibility is that correct mathematical derivations have intrinsic structural regularities (simpler numbers, consistent sign patterns, shorter expressions) that make them marginally more compressible even when the model has never seen them. Another possibility is format or length artifacts. We present J5 as an interesting observation requiring further controls, not as an established result.

\section{Denoising Experiment Per-Seed Details}
\label{app:seeds}

\begin{table}[t]
\caption{J1 per-seed accuracy.}
\label{tab:j1_seeds}
\begin{center}
\begin{tabular}{lccccc}
\toprule
Size & Seed 1 & Seed 2 & Seed 3 & Seed 4 & Mean \\
\midrule
tiny & 64.2\% & 64.1\% & 67.0\% & 65.8\% & 65.3\% \\
small & 75.1\% & 76.3\% & 72.6\% & 74.2\% & 74.6\% \\
medium & 79.6\% & 80.9\% & 82.4\% & 81.3\% & 81.1\% \\
large & 83.5\% & 86.8\% & -- & -- & 85.2\% \\
\bottomrule
\end{tabular}
\end{center}
\end{table}

\begin{table}[t]
\caption{J2 per-seed accuracy.}
\label{tab:j2_seeds}
\begin{center}
\begin{tabular}{lccccc}
\toprule
Size & Seed 1 & Seed 2 & Seed 3 & Seed 4 & Mean \\
\midrule
tiny & 40.2\% & 44.3\% & 46.3\% & 43.0\% & 43.5\% \\
small & 44.8\% & 45.3\% & 40.3\% & 47.4\% & 44.5\% \\
medium & 45.8\% & 40.9\% & 47.7\% & 48.6\% & 45.8\% \\
large & 51.6\% & 50.4\% & -- & -- & 51.0\% \\
\bottomrule
\end{tabular}
\end{center}
\end{table}

\begin{table}[t]
\caption{J3--J5 per-seed accuracy.}
\label{tab:j345_seeds}
\begin{center}
\begin{tabular}{llcccc}
\toprule
Condition & Size & Seed 1 & Seed 2 & Mean \\
\midrule
J3 & tiny & 58.4\% & 59.5\% & 59.0\% \\
J3 & small & 67.5\% & 69.8\% & 68.6\% \\
J3 & medium & 73.4\% & 73.7\% & 73.6\% \\
J3 & large & 74.7\% & 75.7\% & 75.2\% \\
\midrule
J4 & tiny & 57.3\% & 56.0\% & 56.6\% \\
J4 & small & 64.1\% & 65.7\% & 64.9\% \\
J4 & medium & 67.0\% & 66.2\% & 66.6\% \\
J4 & large & 66.1\% & 65.5\% & 65.8\% \\
\midrule
J5 & tiny & 54.1\% & 53.4\% & 53.8\% \\
J5 & small & 60.6\% & 59.2\% & 59.9\% \\
J5 & medium & 60.0\% & 56.2\% & 58.1\% \\
J5 & large & 54.0\% & 58.1\% & 56.1\% \\
\bottomrule
\end{tabular}
\end{center}
\end{table}

\end{document}
